% GENERATED FILE. Edit canonical lessons, then run npm run generate:books.
% canonical-source-hash: fnv1a64:9f062cb2c82d5cf2
% canonical-lessons: MR-C353-pusti-karne, MR-C353-dujora-dene, MR-C353-khandan-karne, MR-C353-vastavik, MR-C353-pratikriya-dene

\chapter{Confirming and Refuting}
\label{ch:mr-confirming-and-refuting}
\hlchaptermodality{353}

\begin{chapteropening}
\textbf{By the end of this chapter:} \emph{I can confirm a report or argue against a claim.}

Complete the last lesson of chapter 353: I can confirm a report or argue against a claim.
\end{chapteropening}

\section[\texorpdfstring{\mr{पुष्टी} \mr{करणे} (puṣṭī karṇe)}{पुष्टी करणे (puṣṭī karṇe)}]{\mr{पुष्टी} \mr{करणे} (puṣṭī karṇe) — to confirm}

\label{lesson:MR-C353-pusti-karne}



\begin{quote}
Before the new one: say the Marathi for especially, then the Marathi for leading.
\end{quote}

\subsection*{You'll want to know: \mr{पुष्टी} \mr{करणे}}
\textbf{\mr{पुष्टी} \mr{करणे}} — \emph{puṣṭī karṇe} — \textquotedblleft{}to confirm\textquotedblright{}.

\begin{quote}
\mr{पोलिसांनी} \mr{अपघाताची} \mr{पुष्टी} \mr{केली}.
\end{quote}

\begin{quote}
\emph{polisānnī apghātācī puṣṭī kelī.}
\end{quote}

\begin{quote}
The police confirmed the accident.
\end{quote}

Past, with -\mr{ने} (he …): \mr{त्याने} \mr{पुष्टी} \mr{केली}, \emph{tyāne puṣṭī kelī}.

\mr{पुष्टी} is feminine, so the past is \mr{केली}.

\begin{cousinweb}
\textbf{\mr{पुष्टी} \mr{करणे}} (\emph{puṣṭī karṇe}) joins the tatsama \mr{पुष्टी}, Sanskrit \emph{puṣṭi}, \textquotedblleft{}growth, thriving\textquotedblright{} (Monier-Williams), to \mr{करणे}, \textquotedblleft{}to make\textquotedblright{}: to strengthen a report until it is sure.
\end{cousinweb}

\begin{grammarlens}[title={What you've built}]
One more word for giving reasons and opinions.
\end{grammarlens}

\subsection*{Your turn}
\begin{itemize}
\raggedright
  \item \emph{Say it:} \emph{puṣṭī karṇe}
  \item \emph{Say it:} \emph{puṣṭī karṇe}, once more
  \item \emph{Say it:} \emph{polisānnī apghātācī puṣṭī kelī.} — The police confirmed the accident.
  \item \emph{Recall:} say the Marathi for especially, then the Marathi for leading, then say \emph{puṣṭī karṇe} again
\end{itemize}

\begin{tcolorbox}[breakable,colback=teal!4,colframe=teal!35!black,title={Before you move on}]
What does \mr{पुष्टी} \mr{करणे} mean? (\textquotedblleft{}To confirm\textquotedblright{}.) Say it once more.
\end{tcolorbox}
\section[\texorpdfstring{\mr{दुजोरा} \mr{देणे} (dujorā deṇe)}{दुजोरा देणे (dujorā deṇe)}]{\mr{दुजोरा} \mr{देणे} (dujorā deṇe) — to corroborate}

\label{lesson:MR-C353-dujora-dene}



\begin{quote}
Before the new one: say the Marathi for leading, then the Marathi for to confirm.
\end{quote}

\subsection*{You'll want to know: \mr{दुजोरा} \mr{देणे}}
\textbf{\mr{दुजोरा} \mr{देणे}} — \emph{dujorā deṇe} — \textquotedblleft{}to corroborate\textquotedblright{}.

\begin{quote}
\mr{मीराच्या} \mr{बोलण्याला} \mr{तिच्या} \mr{भावाने} \mr{दुजोरा} \mr{दिला}.
\end{quote}

\begin{quote}
\emph{mīrācyā bolaṇyālā ticyā bhāvāne dujorā dilā.}
\end{quote}

\begin{quote}
Her brother backed up what Meera said.
\end{quote}

Past, with -\mr{ने} (he …): \mr{त्याने} \mr{दुजोरा} \mr{दिला}, \emph{tyāne dujorā dilā}.

The statement backed takes -\mr{ला}.

\begin{cousinweb}
\textbf{\mr{दुजोरा} \mr{देणे}} (\emph{dujorā deṇe}) joins \mr{दुजोरा}, \textquotedblleft{}support, backing\textquotedblright{}, a Marathi noun whose origin is not settled, to \mr{देणे}, \textquotedblleft{}to give\textquotedblright{}.
\end{cousinweb}

\begin{grammarlens}[title={What you've built}]
One more word for giving reasons and opinions.
\end{grammarlens}

\subsection*{Your turn}
\begin{itemize}
\raggedright
  \item \emph{Say it:} \emph{dujorā deṇe}
  \item \emph{Say it:} \emph{dujorā deṇe}, once more
  \item \emph{Say it:} \emph{mīrācyā bolaṇyālā ticyā bhāvāne dujorā dilā.} — Her brother backed up what Meera said.
  \item \emph{Recall:} say the Marathi for leading, then the Marathi for to confirm, then say \emph{dujorā deṇe} again
\end{itemize}

\begin{tcolorbox}[breakable,colback=teal!4,colframe=teal!35!black,title={Before you move on}]
What does \mr{दुजोरा} \mr{देणे} mean? (\textquotedblleft{}To corroborate\textquotedblright{}.) Say it once more.
\end{tcolorbox}
\section[\texorpdfstring{\mr{खंडन} \mr{करणे} (khaṇḍan karṇe)}{खंडन करणे (khaṇḍan karṇe)}]{\mr{खंडन} \mr{करणे} (khaṇḍan karṇe) — to refute}

\label{lesson:MR-C353-khandan-karne}



\begin{quote}
Before the new one: say the Marathi for to confirm, then the Marathi for to corroborate.
\end{quote}

\subsection*{You'll want to know: \mr{खंडन} \mr{करणे}}
\textbf{\mr{खंडन} \mr{करणे}} — \emph{khaṇḍan karṇe} — \textquotedblleft{}to refute\textquotedblright{}.

\begin{quote}
\mr{त्याने} \mr{एकेक} \mr{मुद्दा} \mr{घेऊन} \mr{आरोपांचं} \mr{खंडन} \mr{केलं}.
\end{quote}

\begin{quote}
\emph{tyāne ekek muddā gheūn āropāñcaṁ khaṇḍan kelaṁ.}
\end{quote}

\begin{quote}
He refuted the accusations point by point.
\end{quote}

Past, with -\mr{ने} (he …): \mr{त्याने} \mr{खंडन} \mr{केलं}, \emph{tyāne khaṇḍan kelaṁ}.

Formal; the claim refuted takes -\mr{चं}: \mr{आरोपांचं} \mr{खंडन}.

\begin{cousinweb}
\textbf{\mr{खंडन} \mr{करणे}} (\emph{khaṇḍan karṇe}) joins the tatsama \mr{खंडन}, Sanskrit \emph{khaṇḍana}, \textquotedblleft{}breaking\textquotedblright{} (CDIAL 3796), to \mr{करणे}, \textquotedblleft{}to do\textquotedblright{}: to break an argument.
\end{cousinweb}

\begin{grammarlens}[title={What you've built}]
One more word for giving reasons and opinions.
\end{grammarlens}

\subsection*{Your turn}
\begin{itemize}
\raggedright
  \item \emph{Say it:} \emph{khaṇḍan karṇe}
  \item \emph{Say it:} \emph{khaṇḍan karṇe}, once more
  \item \emph{Say it:} \emph{tyāne ekek muddā gheūn āropāñcaṁ khaṇḍan kelaṁ.} — He refuted the accusations point by point.
  \item \emph{Recall:} say the Marathi for to confirm, then the Marathi for to corroborate, then say \emph{khaṇḍan karṇe} again
\end{itemize}

\begin{tcolorbox}[breakable,colback=teal!4,colframe=teal!35!black,title={Before you move on}]
What does \mr{खंडन} \mr{करणे} mean? (\textquotedblleft{}To refute\textquotedblright{}.) Say it once more.
\end{tcolorbox}
\section[\texorpdfstring{\mr{वास्तविक} (vāstavik)}{वास्तविक (vāstavik)}]{\mr{वास्तविक} (vāstavik) — actual; real}

\label{lesson:MR-C353-vastavik}



\begin{quote}
Before the new one: say the Marathi for to corroborate, then the Marathi for to refute.
\end{quote}

\subsection*{You'll want to know: \mr{वास्तविक}}
\textbf{\mr{वास्तविक}} — \emph{vāstavik} — \textquotedblleft{}actual; real\textquotedblright{}.

\begin{quote}
\mr{त्याने} \mr{सांगितलेलं} \mr{कारण} \mr{वास्तविक} \mr{नव्हतं}.
\end{quote}

\begin{quote}
\emph{tyāne sāṅgitlelaṁ kāraṇ vāstavik navhataṁ.}
\end{quote}

\begin{quote}
The reason he gave was not the real one.
\end{quote}

Also in fact: \mr{वास्तविक} \mr{पाहता}, in actual fact.

\begin{cousinweb}
\textbf{\mr{वास्तविक}} (\emph{vāstavik}) is a tatsama, taken from Sanskrit \emph{vāstavika}, \textquotedblleft{}substantial, real\textquotedblright{} (Monier-Williams).
\end{cousinweb}

\begin{grammarlens}[title={What you've built}]
One more word for giving reasons and opinions.
\end{grammarlens}

\subsection*{Your turn}
\begin{itemize}
\raggedright
  \item \emph{Say it:} \emph{vāstavik}
  \item \emph{Say it:} \emph{vāstavik}, once more
  \item \emph{Say it:} \emph{tyāne sāṅgitlelaṁ kāraṇ vāstavik navhataṁ.} — The reason he gave was not the real one.
  \item \emph{Recall:} say the Marathi for to corroborate, then the Marathi for to refute, then say \emph{vāstavik} again
\end{itemize}

\begin{tcolorbox}[breakable,colback=teal!4,colframe=teal!35!black,title={Before you move on}]
What does \mr{वास्तविक} mean? (\textquotedblleft{}Actual; real\textquotedblright{}.) Say it once more.
\end{tcolorbox}
\section[\texorpdfstring{\mr{प्रतिक्रिया} \mr{देणे} (pratikriyā deṇe)}{प्रतिक्रिया देणे (pratikriyā deṇe)}]{\mr{प्रतिक्रिया} \mr{देणे} (pratikriyā deṇe) — to react; to respond}

\label{lesson:MR-C353-pratikriya-dene}



\begin{quote}
Before the new one: say the Marathi for to refute, then the Marathi for actual.
\end{quote}

\subsection*{You'll want to know: \mr{प्रतिक्रिया} \mr{देणे}}
\textbf{\mr{प्रतिक्रिया} \mr{देणे}} — \emph{pratikriyā deṇe} — \textquotedblleft{}to react; to respond\textquotedblright{}.

\begin{quote}
\mr{या} \mr{बातमीवर} \mr{तिने} \mr{लगेच} \mr{प्रतिक्रिया} \mr{दिली}.
\end{quote}

\begin{quote}
\emph{yā bātmīvar tine lagec pratikriyā dilī.}
\end{quote}

\begin{quote}
She reacted to the news at once.
\end{quote}

Past, with -\mr{ने} (he …): \mr{त्याने} \mr{प्रतिक्रिया} \mr{दिली}, \emph{tyāne pratikriyā dilī}.

What is reacted to takes -\mr{वर}.

\begin{cousinweb}
\textbf{\mr{प्रतिक्रिया} \mr{देणे}} (\emph{pratikriyā deṇe}) joins the tatsama \mr{प्रतिक्रिया}, Sanskrit \emph{pratikriyā}, \textquotedblleft{}requital, counteraction\textquotedblright{} (Monier-Williams), to \mr{देणे}, \textquotedblleft{}to give\textquotedblright{}.
\end{cousinweb}

\begin{grammarlens}[title={What you've built}]
One more word for giving reasons and opinions.
\end{grammarlens}

\subsection*{Your turn}
\begin{itemize}
\raggedright
  \item \emph{Say it:} \emph{pratikriyā deṇe}
  \item \emph{Say it:} \emph{pratikriyā deṇe}, once more
  \item \emph{Say it:} \emph{yā bātmīvar tine lagec pratikriyā dilī.} — She reacted to the news at once.
  \item \emph{Recall:} say the Marathi for to refute, then the Marathi for actual, then say \emph{pratikriyā deṇe} again
\end{itemize}

\begin{tcolorbox}[breakable,colback=teal!4,colframe=teal!35!black,title={Before you move on}]
What does \mr{प्रतिक्रिया} \mr{देणे} mean? (\textquotedblleft{}To react; to respond\textquotedblright{}.) Say it once more.
\end{tcolorbox}
